\chapter{The data model}

\section{Seven tables}

The database is small on purpose. Seven tables cover the whole application.

\begin{table}[h]
\centering
\begin{tabular}{@{}p{4.4cm}p{9.3cm}@{}}
\toprule
\textbf{Table} & \textbf{Holds} \\
\midrule
\code{users} & The single account that owns the data. \\
\code{categories} & Optional catalogue of category names and colours. \\
\code{transactions} & One expense or one income entry. \\
\code{broker\_imports} & One row per uploaded statement. \\
\code{investments} & A valued position, optionally linked to its import. \\
\code{bank\_accounts} & An account, a savings pot or cash. \\
\code{net\_worth\_snapshots} & Net worth on a given day. \\
\bottomrule
\end{tabular}
\caption{Tables and what they represent.}
\end{table}

There is no table of users' preferences, no audit log and no soft deletes. The
application serves one person; features that only pay off with many users would
cost complexity and return nothing.

\section{Money is never a float}

Two type aliases are declared once and reused in every model:

\begin{lstlisting}[style=py]
MONEY = Numeric(14, 2)
QTY = Numeric(20, 8)
\end{lstlisting}

\code{MONEY} covers amounts: fourteen digits with two decimals. \code{QTY}
covers quantities, which need far more decimals because a crypto position can be
a fraction of a coin.

The reason is the classic one. In binary floating point, \code{0.1 + 0.2} is not
\code{0.3}. Adding a few hundred expenses that way produces a total that is off
by cents, and a total that is off by cents cannot be checked against a bank
statement. Python's \code{Decimal} and SQL's \code{NUMERIC} do exact decimal
arithmetic, so the total either matches or the data is wrong; the arithmetic
itself is never in doubt.

\begin{keypoint}[The rule in practice]
Parsers return \code{Decimal}. The ORM stores \code{NUMERIC}. The backup writes
decimals as strings, not JSON numbers. Floats appear only in the payloads sent
to the charting library, where a rounding error is smaller than a pixel.
\end{keypoint}

\section{Transactions}

\begin{lstlisting}[style=py]
class Transaction(Base):
    __tablename__ = "transactions"

    id: Mapped[int] = mapped_column(primary_key=True)
    date: Mapped[dt.date] = mapped_column(Date, index=True)
    year: Mapped[int] = mapped_column(Integer, index=True)
    month: Mapped[int] = mapped_column(Integer, index=True)
    category: Mapped[str] = mapped_column(String(120), index=True)
    kind: Mapped[str] = mapped_column(String(10))       # expense | income
    amount: Mapped[Decimal] = mapped_column(MONEY)
    note: Mapped[str | None] = mapped_column(String(255), nullable=True)
\end{lstlisting}

Two choices in that model deserve an explanation.

\subsection{Why year and month are stored separately}

The date is already there, so storing the year and the month again looks like
duplication. It is, and it is deliberate.

Every view in the application groups by month. Extracting a month from a date
inside a query means calling a database function, and those functions differ
between Postgres and SQLite; worse, an index on the date column usually cannot
be used once the value is wrapped in a function call. Two indexed integer columns
behave identically on both engines and let the planner use the index.

The cost is that the two representations could disagree. They cannot in practice,
because every write goes through a handful of functions that set all three
fields together.

\subsection{Why the category is a string}

There is a \code{categories} table, but transactions store the category name
directly instead of a foreign key. That means a user can type a new category and
it simply exists, with no extra step and no dialogue asking to create it first.

The price is that renaming a category later means an update across rows. For a
single-user application with a few thousand transactions, that trade is clearly
worth it.

\section{Investments}

\begin{lstlisting}[style=py]
class Investment(Base):
    __tablename__ = "investments"

    broker: Mapped[str] = mapped_column(String(30), index=True)
    asset: Mapped[str] = mapped_column(String(200))
    isin: Mapped[str | None]
    quantity: Mapped[Decimal | None] = mapped_column(QTY, nullable=True)
    invested: Mapped[Decimal] = mapped_column(MONEY, default=Decimal("0"))
    current_value: Mapped[Decimal] = mapped_column(MONEY, default=Decimal("0"))
    withdrawn: Mapped[Decimal] = mapped_column(MONEY, default=Decimal("0"))
    profit: Mapped[Decimal] = mapped_column(MONEY, default=Decimal("0"))
    valued_on: Mapped[dt.date] = mapped_column(Date, index=True)

    # automatic revaluation
    yahoo_symbol: Mapped[str | None]
    current_price: Mapped[Decimal | None] = mapped_column(QTY, nullable=True)
    price_updated_at: Mapped[dt.datetime | None]
    auto_value: Mapped[bool] = mapped_column(Boolean, default=True)

    monthly_contribution: Mapped[Decimal] = mapped_column(MONEY, default=Decimal("0"))
\end{lstlisting}

A row is a position \emph{as valued on a date}, not a position in general. Every
statement import writes a fresh set of rows with today's \code{valued\_on}, and
the old rows stay. That gives a free history of the portfolio, and it means an
import can never destroy earlier data.

The consequence is that reading "the portfolio" means reading, for each broker,
only the rows of its most recent valuation. That query lives in one function,
\code{latest\_investments()}, and everything else calls it.

\begin{lstlisting}[style=py]
def latest_investments(db: Session) -> list[Investment]:
    rows = []
    for broker in db.scalars(select(Investment.broker).distinct()).all():
        max_date = db.scalar(
            select(func.max(Investment.valued_on)).where(Investment.broker == broker)
        )
        rows.extend(db.scalars(select(Investment).where(
            Investment.broker == broker, Investment.valued_on == max_date)).all())
    return rows
\end{lstlisting}

The loop runs per broker rather than once for the whole table because brokers
are imported on different days. Taking a single global maximum date would hide
every broker that was not imported most recently.

\section{Bank accounts and net worth snapshots}

Bank accounts are deliberately dumb: a name, a kind, a balance and a timestamp.
The user updates the balance when convenient; the application never tries to
connect to a bank.

\code{net\_worth\_snapshots} holds one row per day with three amounts:
investments, banks and the total. It is written by an upsert, so several changes
on the same day update the same row.

\begin{lstlisting}[style=py]
def snapshot_today(db: Session) -> NetWorthSnapshot:
    inv_value, bank_value = _net_worth_parts(db)
    today = date.today()
    snap = db.scalar(select(NetWorthSnapshot).where(NetWorthSnapshot.date == today))
    if snap is None:
        snap = NetWorthSnapshot(date=today)
        db.add(snap)
    snap.investments = Decimal(str(round(inv_value, 2)))
    snap.banks = Decimal(str(round(bank_value, 2)))
    snap.total = snap.investments + snap.banks
    db.commit()
    return snap
\end{lstlisting}

\begin{keypoint}[Why store history instead of computing it]
Net worth could be recomputed from current values whenever the chart is drawn.
It would also redraw the past every time a price moved, because the only prices
available are today's. Storing a snapshot freezes what was true on the day, which
is the whole point of an evolution chart.
\end{keypoint}

The snapshot is written after any change that can move net worth: editing an
investment, contributing to one, editing an account, or the nightly job. Missing
one of those calls does not corrupt anything, it only leaves a gap in the chart.

\section{Migrations}

Schema changes go through Alembic. The environment file reads the URL and the
metadata from the application itself, so there is no second copy of the
connection settings:

\begin{lstlisting}[style=py]
config.set_main_option("sqlalchemy.url", settings.sqlalchemy_url)
target_metadata = Base.metadata
...
context.configure(connection=connection, target_metadata=target_metadata,
                  render_as_batch=settings.is_sqlite)
\end{lstlisting}

\code{render\_as\_batch} matters when developing against SQLite, which cannot
alter most things in place. In batch mode Alembic creates a new table, copies the
data and swaps it, which reproduces on SQLite what Postgres does with a plain
\code{ALTER TABLE}.

The usual cycle is two commands:

\begin{lstlisting}[style=sh]
alembic revision --autogenerate -m "what changed"
alembic upgrade head
\end{lstlisting}

Autogenerate compares the models with the live database and writes the
difference. It is a starting point, not an oracle: it does not see renames, and
it does not know how to move data. Every generated file is worth reading before
it is applied.
